% Appendix A -- Worksheet Capacity (in place of the model's "Managing
% Memory", annex Sec 6 App. C).  Sources (Law 1): os/i123/app/face_a.c
% (SHEET_CELLS 2048, the 32K entry pool and the 4K record area of SHEET_MEM;
% the 64K file buffer g_wk; the typed-entry limit SCR_COLS - 2 = 74; the
% 48-character file name g_name; the prompt ranges 1..240, 0..15, 1..50;
% the two-bank transaction of every command), os/i123/core/i123.h
% (I123_MAX_COLS/ROWS, I123_CODE_MAX), spec/i123/era_shortfalls.md
% ("Capacity is an explicit implementation limit"), ADR-0009 Amendment
% DEC-01a (Initech 123 computes with the 80387-class coprocessor).
% GATES: test-i123-wk1-roundtrip test-i123-quit test-flair-i123-files test-pubs-i123
\Gates{test-i123-wk1-roundtrip, test-i123-quit, test-flair-i123-files, test-pubs-i123}
\ChapterTab{Capacity}
\chapter{Worksheet Capacity}

This appendix lists the limits of Initech 123 and explains what happens
when a worksheet reaches them.

\section{The Limits}
\begin{LotTable}{p{60mm}p{68mm}}
\LotHead{\textbf{Item} & \textbf{Limit}}
Columns & 256 (A through IV)\\
Rows & 8,192\\
Cells with entries or formats & 2,048\\
Characters in labels and formulas, all cells together & about 32,000\\
Range names and other worksheet records & about 4,000 characters\\
Characters in an entry you type & 74\\
Characters in a range name & 15\\
Column width & 1 to 240 characters\\
Decimal places of a format & 0 to 15\\
Iterations of a recalculation & 1 to 50\\
Size of a worksheet file Initech 123 can retrieve & 64K (65,536 bytes)\\
Characters in a file name typed at a prompt & 48\LotEnd
\end{LotTable}

A cell counts toward the limit of 2,048 when it holds an entry, and also when
it is empty but has a format or protection setting of its own, as a cell
does after \slashcmd{/Range Erase} or \slashcmd{/Range Format}.

\section{When the Worksheet Is Full}
Every command that changes the worksheet builds the changed worksheet
first and puts it in place only when it is complete.  When a command would
need more cells or more space than the limits allow, Initech 123 displays
\Msg{Not enough worksheet memory}, and the worksheet is left exactly as it
was before the command.  The same message appears, and the worksheet in
memory is kept, when you retrieve a file that holds more entries than
Initech 123 can take.

An entry you type that does not fit is refused: Initech 123 switches to
\ind{EDIT} mode and leaves the entry on the control panel.  Press
\kname{ESC} to abandon it.

If the worksheet in memory cannot be written as a file
(\Msg{Worksheet too large to save}), /File Save leaves the file on disk
untouched.  Erase entries you do not need, or move part of the work to a
second worksheet, and save again.

\section{Numbers}
Initech 123 stores every number in the IEEE double format of the 80387 and
80486 numeric coprocessor, with about 15 significant digits, and calculates
with the coprocessor.  The General format displays as many of those digits
as the column allows.
